A held-out score on one corpus said little about how the flow detectors we tested behave on traffic from another network. In 5G-NIDD, \RCopiesPct\% of the benign flows are the UDP flood of the other base station, recorded packet for packet in station~1's capture, and a published 99.9\% accuracy rests on record-position fields that a deployed detector never sees. Once the labels and the byte accounting of the exporters are corrected, tree-based detectors trained on NetsLab-5GORAN-IDD score as well on 5G-NIDD as on held-out source data, but \RGapTreeCleanMin{} to \RGapTreeCleanMax{} below references trained on 5G-NIDD and tested on the same flows; on the 5G core testbed they catch the flood they were trained on and flag much benign traffic. At a declared prevalence of 0.002 and bounds of our own, no detector passed our two-term test. The best detectors rank the attacks of 5G-NIDD well enough for the bounds, but a threshold fixed on source data did not carry over, and none of three adaptation methods without target labels closed the gap, whereas one labeled target capture closed most of it. On the radio layer, the random-split gain was not significant under either grouping, and the latency we measured for window-level detectors in a FlexRIC loop with an emulated E2 node (99th percentile at most \RicLoopTenPnnMax{}\,ms) is a processing cost on one host, not a conformance result; flow detectors were not timed. We recommend that studies of 5G flow detectors report at least:
\begin{enumerate}
\item performance on an independently collected corpus in a prevalence-invariant metric, beside input-blind baselines and a reference trained on that corpus and tested on the same flows;
\item the split protocol and group key, and intervals that account for overlapping splits and clustered samples;
\item how each flow exporter counts bytes, and how the shared feature space removes the differences;
\item evidence that the labels are a function of the features used, such as the share of identical records under opposite labels, and that no record-position or identifier field is a feature;
\item the false positive rate with its unit, the alert rate and operational precision at a stated base rate, and whether the operating threshold was chosen with target labels;
\item for a detector meant for the RIC, the latency from indication to control acknowledgment, with stages, implementation, and host, measured with a real E2 node before any conformance claim.
\end{enumerate}
